% Generated draft: requires mathematical and source review.
\documentclass[11pt]{article}
\usepackage[margin=25mm]{geometry}
\usepackage{fontspec}
\setmainfont{Latin Modern Roman}
\usepackage{amsmath,amssymb,mathtools,unicode-math}
\setmathfont{Latin Modern Math}
\usepackage{enumitem,needspace,xurl,hyperref}
\setlength{\parindent}{0pt}
\setlength{\parskip}{5pt}
\setlength{\emergencystretch}{3em}
\newcommand{\N}{\mathbb{N}}
\newcommand{\Z}{\mathbb{Z}}
\newcommand{\Q}{\mathbb{Q}}
\newcommand{\R}{\mathbb{R}}
\newcommand{\C}{\mathbb{C}}
\newcommand{\F}{\mathbb{F}}
\newcommand{\A}{\mathbb{A}}
\newcommand{\PP}{\mathbb{P}}
\newcommand{\E}{\mathbb{E}}
\newcommand{\eps}{\varepsilon}
\newcommand{\dd}{\,\mathrm{d}}
\begin{document}
\section*{UnsolvedMath ID: 30000311}
\textbf{Source problem code:} catalogue entry
\subsection{Definitions and mathematical statement}
\textbf{Construction request.} Explain how a multiplicative unitary gives rise to a pair of Hopf \(C^*\)-algebras. This is phrased as an explanatory construction request, not by itself a precise new open conjecture or classification theorem. Let \(\mathcal H\) be a Hilbert space. Under one standard convention a multiplicative unitary is a unitary \(W\in\mathcal B(\mathcal H\otimes\mathcal H)\) satisfying the pentagon equation\[W_{12}W_{13}W_{23}=W_{23}W_{12}\quad\text{on }\mathcal H^{\otimes3},\] where subscripts specify the tensor legs. Other conventions reverse or flip this equation and change the formulas below accordingly.\par The customary candidate legs are\[A=\overline{\operatorname{span}}^{\|\cdot\|}\{(\operatorname{id}\otimes\omega)(W):\omega\in\mathcal B(\mathcal H)_*\},\qquad \widehat A=\overline{\operatorname{span}}^{\|\cdot\|}\{(\omega\otimes\operatorname{id})(W):\omega\in\mathcal B(\mathcal H)_*\},\] where \(\mathcal B(\mathcal H)_*\) denotes the normal linear functionals. With this convention the intended coproduct formulas are\[\Delta(a)=W^*(1\otimes a)W,\qquad \widehat\Delta(b)=\Sigma W(b\otimes1)W^*\Sigma,\] where \(\Sigma(\xi\otimes\eta)=\eta\otimes\xi\). The construction must state hypotheses ensuring that the slice spaces are nondegenerately represented \(C^*\)-algebras and that these formulas define nondegenerate *-homomorphisms into \(M(A\otimes_{\min}A)\) and \(M(\widehat A\otimes_{\min}\widehat A)\), respectively. Here \(M(C)\) is the multiplier algebra. One must then establish coassociativity, \((\Delta\otimes\operatorname{id})\Delta=(\operatorname{id}\otimes\Delta)\Delta\), and the dual analogue, together with whichever cancellation, counit, antipode, or Haar-weight properties the intended definition of Hopf \(C^*\)-algebra requires. For example, a cancellation convention asks that the closed linear spans of \(\Delta(A)(1\otimes A)\) and \((A\otimes1)\Delta(A)\) equal \(A\otimes_{\min}A\).\par The source does not specify regularity, manageability, or another sufficient class of multiplicative unitaries, and the bare pentagon equation is not silently asserted sufficient for every analytic property just listed. Nor is a pseudo-multiplicative unitary over a nontrivial base identified with the ordinary tensor-product object above. These missing hypotheses and the meaning of Hopf \(C^*\)-algebra are essential to a rigorous version of the requested explanation.\par The supplied triage checked 2026-08-21 labels the item partially solved and says multiplicative and pseudo-multiplicative unitary theory yields substantial Hopf-algebraic constructions while a broad correspondence is not completely classified. Such a classification is broader than the wording of the original explanatory request. The dated assessment is retained without turning it into a specific unresolved theorem or an independently checked literature conclusion.
\subsection{Short English statement}
Describe the two slice algebras and their coproducts associated with a multiplicative unitary, under hypotheses that make the construction valid. The pentagon convention and analytic conditions must be stated. The source asks for an explanation; its broader classification-status note is not itself a precise conjecture.
\subsection{Sources}
\begin{itemize}\item Local UnsolvedMath registry, permanent ID 30000311, code catalogue entry; \emph{Hopf \(C^*\)-Algebras from Multiplicative Unitaries}, Oberwolfach Reports Open Problems.\item \emph{C*-Algebren (2006)}, Oberwolfach Reports, DOI 10.4171/owr/2005/41, pp. 2305--2374. \url{https://doi.org/10.4171/owr/2005/41}. Proposed year and year tag: 2005; no proposer named.\item Supplied literature label \emph{original Oberwolfach report}: \url{https://doi.org/10.4171/owr/2005/41}. No separate authoritative construction theorem or precise sufficient-hypothesis citation is supplied.\item Literature triage checked 2026-08-21, status \texttt{partially\_solved}; record updated 2026-09-16. The \texttt{corrected\_verified} label concerns restoration of \(C^*\) notation. Verify the dated assessment independently.\end{itemize}
\end{document}
